% problem-000113 2 双氧物种电子与晶体
\section*{2. 双氧物种电子与晶体}
\textit{下列为分问。}

\subsection*{2-1}
请解释为什么在DNA分⼦中，腺嘌呤与胸腺嘧啶含量相等，⻦嘌呤与胞嘧啶含量相等。

A
C
（图：）
腺嘌呤(A)

（图：）
⻦嘌呤(G)

（图：）
胸腺嘧啶(T)

（图：）
图1 碱基结构
胞嘧啶(C)

\subsection*{2-2}
钾氩定年法可⽤于分析⽕星矿物形成的年代。已知40K衰变时，10.5%为β+衰变⽣成40Ar，89.5 %为β−衰变⽣成 $^ { 4 0 } \mathrm { C a }$ ，40K的半衰期为12.5亿年。试求40K发⽣β+衰变和β−衰变的速率常数 $\cdot k _ { 1 }$ 和 $k _ { 2 }$ 。

\subsection*{2-3}
写出与 $\mathrm { N _ { 2 } H _ { 4 } }$ 质⼦总数相同且含有⾮极性键的物质的分⼦式。

\subsection*{2-4}
画出分⼦式为 $\mathrm { C _ { 3 } O _ { 3 } C l _ { 6 } }$ 中所有Cl原⼦等价的物质的结构式。

\subsection*{2-5}
⼀般化学交联的聚合物难以进⾏⼆次利⽤，但如图2所⽰的A、B、C三种单体共聚后经紫外光照射可以形成交联聚合物，且该交联聚合物能够实现回收热塑加 $\scriptstyle \sum _ { 0 }$ 。请简述该交联聚合物可热塑加⼯的基本原理。

（图：）

（图：）

（图：）
图2 三种单体结构
